\begin{abstract}
Parkour became a FIG World-Championship discipline in 2018 and has a
149-entry Code of Points~2025 that closes the door on ad-hoc trick
naming. Despite this, no open-source system exists that takes a
monocular competition video and produces a FIG-compliant D-score.
We present \sys{}, a deliberately modular pipeline that combines
an inversion-based trick segmenter, an off-the-shelf vision--language
model (VLM), a five-level fuzzy matcher that snaps free-form model
output onto the 149 FIG entries, and a deterministic D-score assembler.
Every component has a text-level interface and can be swapped in
isolation.

We report an honest empirical study of the system on real smartphone
and GoPro parkour footage. Across five trained-model approaches we
abandoned in eighteen months of work (2D pose plus physics, 3D human
mesh recovery with GVHMR, metric learning on 1\,618 parkourtheory
clips, hierarchical attribute classifiers, and CLIP zero-shot), none
reached usable accuracy on competition video. We therefore turned to
frontier VLMs. On a curated 13-clip single-trick benchmark, Claude
Sonnet~4.5 (via composite frame strips), Gemini~2.5, and Qwen2.5-VL
each identify the correct FIG trick on only a small fraction of
clips in headless automation mode---even though the same models
routinely succeed when queried interactively on the same clips.
We document this automation gap, catalogue the recurring failure
modes (twist under-counting, takeoff confusion, wall-context loss),
show that the fuzzy matcher and inversion segmenter remain useful
regardless of which VLM is plugged in, and release the full system,
prompts, ground-truth labels, and raw result files so the community
can build on the components that do work.
\end{abstract}
